\section{Sticking: interlacement is causal past-closure}

This subgraph maps textile interlacement onto chromatic-Lorentzian closure.  The
proved carrier supplies \texttt{self\_le\_settle},
\texttt{box\_pastSum}, and \texttt{settle\_iff\_pastSum\_pos}; the incidence
claim is that crossings, cloth, and anchoring are respectively the settle-step,
retarded propagator, and past-source predicate.

\begin{theorem}[Interlacement is a settle-step]\label{thm:interlacement-is-settle}\lean{Obsmetrologia.AperiodicOrder.Tiling.ChromaticLorentzian.self_le_settle}\leanok\uses{thm:weave-draft-sofic, thm:chromatic-lorentzian-carrier}
Along one warp thread, over/under absorption is witnessed by
\texttt{self\_le\_settle}: $f\,n \lor \texttt{settle}\,f\,n =
\texttt{settle}\,f\,n$.  A float marks the non-absorbed residue.  The open fiber is
the full warp$\times$weft carrier of \cref{thm:weave-2d-sofic}.
\end{theorem}

\begin{theorem}[The woven sheet is the propagator]\label{thm:fabric-is-propagator}\lean{Obsmetrologia.CausalSets.Action.Chromatic.box_pastSum}\leanok\uses{thm:interlacement-is-settle, thm:chromatic-lorentzian-action}
The cloth is the retarded propagator $G$ (\texttt{pastSum}).  The causal
d'Alembertian recovers the local crossing datum:
$\Box\circ G=\mathrm{id}$ by \texttt{box\_pastSum}.  The node identifies weaving
with causal integration and draft-reading with the inverse difference.
\end{theorem}

\begin{theorem}[A thread sticks iff anchored in its past]\label{thm:stick-iff-past-source}\lean{Obsmetrologia.CausalSets.Action.Chromatic.settle_iff_pastSum_pos}\leanok\uses{thm:fabric-is-propagator}
A crossing is locked iff a past source propagates to it:
\texttt{settle\_iff\_pastSum\_pos} gives
$\texttt{settle}\,f\,n \iff 0<\texttt{pastSum}\,f\,n$.  Selvedge anchoring is the
source case; interior sticking is the propagated case.
\end{theorem}
